%% Заключение

В рамках выпускной квалификационной работы была спроектирована и
реализована новая подсистема хранения полуструктурированных данных
для колоночного движка Otterbrix. Проделанная работа охватывает
следующие основные результаты.

Во-первых, проведён обзор существующих в индустрии подходов к
хранению JSON в СУБД: хранение как неделимого значения (BSON,
JSONB в PostgreSQL), модель Entity-Attribute-Value и её вариации,
колоночное развёртывание путей, и гибридный тип \texttt{JSON} в
ClickHouse. Сформулированы требования к новому решению,
вытекающие из недостатков существующих подходов.

Во-вторых, разработана модель гетерогенного хранения с
автоматической миграцией полей между разреженным и плотным
режимами. Часто встречающиеся поля хранятся как полноценные
колонки основной таблицы, редкие~--- в отдельных двухколоночных
вспомогательных таблицах формата \texttt{(\_id, value)}; при
пересечении заранее фиксированного порога заполненности поле
автоматически мигрирует в основную таблицу.

В-третьих, модель реализована в кодовой базе Otterbrix: введён
класс \texttt{computed\_schema}, инкапсулирующий динамическую
схему, расширены пути диспетчера для обработки \textsc{insert} и
\textsc{select}, добавлена поддержка SQL-конструкций
\texttt{CREATE TABLE db.t ();},
\texttt{WITH (sparse\_threshold=\ldots,
pinned\_columns=\ldots)} и
\texttt{INSERT \ldots\ VALUES (json('\ldots'))}. Интеграция с
существующей акторной моделью параллелизма не потребовала введения
новых механизмов синхронизации; согласованность в духе
snapshot~isolation обеспечивается за счёт переиспользования
действующего MVCC-протокола.

В-четвёртых, выполнены оптимизации физических операторов:
проекция колонок на уровне оптимизатора, типизированные предикаты
фильтрации с пакетным копированием подошедших строк,
специализированные пути GROUP~BY, исправление ошибки в реализации
zone~maps, параллельная запись во вспомогательные таблицы. Без
этих оптимизаций теоретически правильная архитектура хранения не
дала бы конкурентоспособной производительности.

В-пятых, проведено сравнительное тестирование на стандартном
бенчмарке JSONBench со сравнением с ClickHouse~24.10 и
PostgreSQL~17. Результаты показали, что предложенная архитектура
достигает производительности, сопоставимой с ClickHouse
(20--30\% разрыв на типичных запросах, до 1.5$\times$ на
специфических сценариях), и существенно (в 10--20~раз) превосходит
PostgreSQL JSONB. Объём хранимых данных близок к показателям
ClickHouse.

Ключевое преимущество предложенного подхода~--- в естественной
интеграции с существующими механизмами колоночного реляционного
движка: разреженные таблицы являются обычными таблицами,
автоматически получают все оптимизации, правильно интегрируются с
MVCC, WAL и индексами. В отличие от ClickHouse, где редкие пути
хранятся в специальном внутреннем формате с ограниченной
функциональностью, предложенное решение сохраняет для таких полей
полный набор операций и оптимизаций. Объём новой кодовой базы
измеряется тысячами строк~--- на порядок меньше того, что
потребовался бы для разработки специализированной подсистемы
хранения с нуля.

Основным направлением дальнейшей работы является перевод
реконструкции \textsc{select} на явные узлы LEFT~JOIN в логическом
плане, что позволит применять стандартные оптимизации оптимизатора
(порядок соединений, push-down предикатов) и обеспечит единообразие
плана запроса в EXPLAIN. Дополнительные направления~--- адаптивный
порог, обратная миграция, сжатие неактивных вспомогательных таблиц,
специализированные индексы для разреженных таблиц, поддержка
распределённой репликации.

Результаты работы служат практическим подтверждением того, что
правильно спроектированная схема хранения, опирающаяся на
существующие реляционные механизмы, способна конкурировать со
специализированными решениями, предъявляя значительно меньшие
требования к объёму новой кодовой базы. Это направление~---
использование существующих реляционных абстракций для хранения
полуструктурированных данных~--- представляется перспективным для
дальнейших исследований.
